% ==============================================================================
% 4. PENUTUP DAN REKOMENDASI PEMBANGUNAN
% ==============================================================================

\section{Penutup dan Rekomendasi Pembangunan}

Berdasarkan kompilasi dan analisis 55 variabel data potensi dari 7 wilayah Jaga yang disajikan dalam publikasi \textbf{``Desa Popontolen Dalam Angka 2026''}, Pemerintah Desa Popontolen merumuskan rekomendasi strategis sebagai arah prioritas kebijakan pembangunan desa:

\begin{enumerate}
  \item \textbf{Pemberdayaan Tenaga Kerja Usia Produktif \& Penguatan UMKM:} \\
  Dengan 942 jiwa (64,4\%) berada pada usia produktif dan \textit{Dependency Ratio} 55,3\%, Desa Popontolen berada pada era bonus demografi. Program pelatihan vokasi, sertifikasi keahlian, dan fasilitasi akses pembiayaan diprioritaskan bagi 34 unit usaha ekonomi (6 minimarket, 18 toko kelontong, 9 warung, 1 restoran).

  \item \textbf{Pengembangan Sentra Pertanian dan Perkebunan Terpadu:} \\
  Sebanyak 179 KK petani (33,6\% dari total keluarga) menjadi basis ketahanan pangan, terkonsentrasi di Jaga 7 (67 KK), Jaga 6 (30 KK), Jaga 2 (27 KK), dan Jaga 4 (24 KK). Diprioritaskan penyaluran bibit kelapa unggul, pupuk bersubsidi, serta hilirisasi produk turunan kelapa.

  \item \textbf{Pemerataan Akses Kelistrikan:} \\
  Terdapat \textbf{2 keluarga di Jaga 6} yang belum memiliki akses listrik dan \textbf{8 keluarga} (Jaga 3: 2 KK, Jaga 4: 6 KK) menggunakan listrik non-PLN. Pemerintah Desa memprioritaskan perluasan jaringan PLN dan program bantuan sambungan listrik gratis.

  \item \textbf{Penguatan Fasilitas Kesehatan dan Inklusi Sosial:} \\
  Dari 13 jenis fasilitas kesehatan yang didata, hanya tersedia \textbf{1 unit Puskesmas Pembantu} di Jaga 6. Terdapat \textbf{23 jiwa penyandang disabilitas} yang memerlukan program jaring pengaman sosial dan pemberdayaan inklusi terpadu.

  \item \textbf{Pengembangan Akses Pendidikan Lanjutan, Perbankan, dan Logistik:} \\
  Desa Popontolen tidak memiliki fasilitas pendidikan SMP dan SMA; tidak ada fasilitas perbankan (6 jenis bank seluruhnya nihil); serta tidak ada kantor pos dan jasa ekspedisi swasta. Pemerintah Desa perlu berkoordinasi dengan pemerintah kabupaten untuk memperluas akses pendidikan lanjutan, inklusi keuangan digital, serta infrastruktur logistik desa.
\end{enumerate}

Demikian publikasi ini disusun dengan harapan dapat memberikan sumbangsih nyata dalam mewujudkan tata kelola desa yang transparan, akuntabel, dan berbasis data.
